\section{Bayesian model comparison: exact proofs}
\label{sec:annex-p15-bayesian-model-comparison}

\begin{annexitem}{annex:p15.model-posterior-odds}{thm:p15-model-posterior-odds}{Integrating the joint model law}
\begin{proof}
By \cref{thm:p15-tonelli}, \(m_k\) is measurable and
\[
 \int m_k\,d\nu=\int_{\Theta_k}\left(\int f_k(x\mid\theta)\nu(dx)\right)
                  \Pi_k(d\theta)=1.
\]
A nonnegative function with finite integral is finite almost everywhere.
The finite union of the exceptional sets for all \(m_k\) is therefore
\(\nu\)-null. The mixture density also integrates to one; its zero set
has zero \(Q_\alpha\)-mass and its infinite set is \(\nu\)-null.
Thus \(Q_\alpha(\mathcal G)=1\).

Construct the prior on the tagged disjoint union
\(\bigsqcup_k(\{k\}\times\Theta_k)\) by assigning
\(\alpha_k\Pi_k(A)\) to \(\{k\}\times A\). A finite tagged union of
standard Borel spaces is standard Borel, and its displayed measure is a
probability. With likelihood \(f_k(x\mid\theta)\), its joint law satisfies
\[
 \Pr\{X\in D,H_k,\theta\in A\}
 =\alpha_k\int_A\int_D f_k(x\mid\theta)\nu(dx)\Pi_k(d\theta).
\]
For \(A=\Theta_k\), Tonelli gives \(\alpha_k\int_Dm_k\,d\nu\).
On \(\mathcal G\), \(w_ke=\alpha_km_k\), so integrating \(w_k\)
against \(Q_\alpha\) over \(D\) gives this model/data joint probability.
The complement of \(\mathcal G\) has zero predictive mass, and the
default vector is a measurable probability vector. This verifies the
conditional model probabilities.

The tagged construction also meets the hypotheses of
\cref{thm:p15-bayes-posterior-kernel}. Where \(m_k(x)>0\), its
conditional parameter law is
\[
 \Pi_{k,x}(A)=\frac{\int_A f_k(x\mid\theta)\Pi_k(d\theta)}{m_k(x)}.
\]
If \(m_k(x)=0\), assign the default \(\Pi_k\); the accompanying
weight is zero on \(\mathcal G\). Multiplication by \(w_k\) gives the
tagged posterior formula, including zero model prior weights. Absolute
continuity of \(\Pi_{k,x}\) with respect to \(\Pi_k\) also shows that
a within-model prior-null parameter event stays null.

If \(\alpha_jm_j(x)>0\), divide the two weights and cancel \(e(x)\).
This yields the odds identity without dividing by a zero prior or likelihood.
If both marginal densities are positive, reciprocal factors multiply to
one, and the logarithm-of-a-positive-quotient identity gives
\(\log B_{ij}=\log m_i-\log m_j\).
The displayed definition of \(B_{ij}\) contains the within-model
priors through \(m_i,m_j\), but not \(\alpha\). For the numerical odds
example, \(20/100=1/5\), and odds \(1/5\) give probability
\((1/5)/(1+1/5)=1/6\).
\end{proof}
\end{annexitem}

\begin{annexitem}{annex:p15.model-density-invariance}{thm:p15-model-density-invariance}{An intrinsic dominating probability and transported priors}
\begin{proof}
Suppose \(f_k,\widetilde f_k\) are jointly measurable versions of the
same conditional data densities for each parameter. For each \(\theta\),
the \(\nu\)-measure of their disagreement set is zero. Apply Tonelli to
the indicator of this measurable disagreement set. It has zero product
measure, and for \(\nu\)-almost every \(x\) the functions agree
\(\Pi_k\)-almost everywhere in \(\theta\). Their integrals are equal
there. Taking the finite intersection over models proves equality of
marginal densities on one common full-measure set. The weight formulas
then agree on a set of full \(Q_\alpha\)-measure.

Put \(\bar m=(K+1)^{-1}\sum_km_k\). This is the density of \(Q_*\).
On the set where \(0<\bar m<\infty\), let \(r_k=m_k/\bar m\).
Set \(r_k=0\) elsewhere. That exceptional set is \(Q_*\)-null:
the infinite-density sets are \(\nu\)-null, and the zero set has zero
\(Q_*\)-mass. On the zero set of \(\bar m\), all \(m_k\) vanish;
therefore for every measurable \(D\),
\[
 \int_D r_k\,dQ_*=\int_Dm_k\,d\nu=Q_k(D).
\]
The ignored infinite sets contribute zero to either integral.
Thus \(r_k\) is a Radon--Nikodym density of \(Q_k\) with respect to
\(Q_*\), and \(0\leq r_k\leq K+1\) almost everywhere.
Uniqueness in \cref{thm:p15-radon-nikodym} makes this description
independent of any alternative common sigma-finite data measure.
Dividing \(m_i/\bar m\) by \(m_j/\bar m\) proves the intrinsic factor
identity where \(r_j>0\). The common factor similarly cancels in model
weights where \(\sum\alpha_kr_k>0\), a set of full \(Q_\alpha\)-mass.

For \(d\nu'=h\,d\nu\), using compatible versions gives
\[
 \int_D(f_k/h)\,d\nu'=\int_Df_k\,d\nu.
\]
The positive finite value \(h(x)\) can be pulled outside the parameter
integral, giving \(m'_k(x)=m_k(x)/h(x)\). It cancels in every defined
ratio. These pointwise equalities require the versions and domain stated
in the theorem; general measure equivalence alone specifies only
almost-everywhere classes.

For a bimeasurable bijection \(T:\Theta_k\to\Theta'_k\), define
\(\Pi'_k(A)=\Pi_k(T^{-1}A)\). For every nonnegative measurable \(g\),
\(\int g\,d\Pi'_k=\int g(T\theta)\Pi_k(d\theta)\): it holds for
indicators by definition, then for nonnegative simple functions and
increasing simple approximations. Apply it to
\(g(\eta)=f_k(x\mid T^{-1}\eta)\) for each \(x\).
This proves exact invariance of the marginal likelihood under the stated
transport, not under an arbitrary reassignment of a coordinate density.

For an observation bijection \(S\), define its pushed-forward measure by
\((S_\#\nu)(D)=\nu(S^{-1}D)\). The same indicator/simple-function
argument proves integration against this pushforward. Hence
\(\int f_k(S^{-1}y\mid\theta)(S_\#\nu)(dy)=1\) for every parameter,
and parameter integration gives \(m'_k(Sx)=m_k(x)\).
Images of a countable finite-measure cover show that \(S_\#\nu\) is
sigma-finite. Thus it is a valid transported data model, with unchanged
ratios at the corresponding comparison points.

Finally, modifying the uniform density at the one-point set \(\{1/2\}\)
does not change any Lebesgue integral. With numerator density changed to
\(c>0\) and denominator density one, the pointwise ratio is \(c\) at
that point and one elsewhere. This verifies the body's null-point example
and why an unqualified pointwise invariance claim would be false.
\end{proof}
\end{annexitem}

\begin{annexitem}{annex:p15.nuisance-integration}{thm:p15-nuisance-integration}{The prior integral and a profile counterexample}
\begin{proof}
For a supplied kernel \(\Lambda(\psi,d\lambda)\) and prior \(\Pi^\psi\),
the disintegration identity first says
\[
 \Pi(A\times B)=\int_A \Lambda(\psi,B)\Pi^\psi(d\psi).
\]
The same integration identity holds for every nonnegative measurable
function of \((\psi,\lambda)\). To justify the extension, it holds for
rectangle indicators. The class of sets for which it holds and whose
section integral is measurable contains the whole space, is closed under
complements (both measures have mass one), and is closed under countable
disjoint unions by additivity and monotone convergence. It is therefore
a Dynkin system containing the rectangle pi-system, so the pi-lambda
theorem gives all product measurable sets. Nonnegative simple functions and monotone convergence
then give the assertion for nonnegative functions. The same rectangle
argument with the unused coordinate \(x\) proves joint measurability of
the kernel integral as a function of \((x,\psi)\).

Apply this identity to \(f_k(x\mid\psi,\lambda)\) for each fixed \(x\).
It gives exactly the iterated marginal formula in the theorem.
A probability-weighted integral of a nonnegative function bounded above
by \(M<\infty\) is at most \(M\); choosing its supremum proves the
upper bound. No supremum measurability or normalization follows from this
inequality.

In the finite example, the two component probabilities at outcomes
\(1,0\) are \((.9,.1)\) and \((.1,.9)\). Their equal-prior average
is \((.5,.5)\), so at 1 the integrated factor against the .6 model is
\(.5/.6=5/6\). Their pointwise maxima are \((.9,.9)\), which sum to
1.8; the maximum-to-null ratio is \(.9/.6=3/2\).
This proves both the reversal and the failed density normalization.
A prior weight \(\beta\) instead gives
\(\beta(.9)+(1-\beta)(.1)=.1+.8\beta\).
\end{proof}
\end{annexitem}

\begin{annexitem}{annex:p15.normal-point-null-factor}{thm:p15-normal-point-null-factor}{Gaussian integration with its constants}
\begin{proof}
Write \(a=s^2>0\), \(b=\tau^2>0\), and \(\mu=\mu_1\).
Completing the square in the product likelihood and prior gives
\[
 \frac{(x-\theta)^2}{a}+\frac{(\theta-\mu)^2}{b}
 =\frac{a+b}{ab}
     \left(\theta-\frac{bx+a\mu}{a+b}\right)^2
   +\frac{(x-\mu)^2}{a+b}.
\]
The product of the two normalizing constants is
\((2\pi\sqrt{ab})^{-1}\). The Gaussian integral in
\cref{thm:p15-normal-normalization} supplies the integral over \(\theta\)
of the first exponential as \(\sqrt{2\pi ab/(a+b)}\).
Consequently
\[
 m_1(x)=\frac1{\sqrt{2\pi(a+b)}}
           \exp\!\left(-\frac{(x-\mu)^2}{2(a+b)}\right).
\]
The point-null prior simply evaluates the likelihood at \(\theta_0\).
Both densities are finite and positive at every real \(x\).
Taking their ratio proves the general expression. If \(\mu=\theta_0\),
substitute \(r=b/a\), \(z=(x-\theta_0)/\sqrt a\); subtraction of the
two quadratic terms gives \(-rz^2/[2(1+r)]\).

For the full sample, expand
\[
 \sum_{i=1}^n(x_i-\theta)^2
 =\sum_{i=1}^n(x_i-\bar x)^2+n(\bar x-\theta)^2,
\]
because \(\sum_i(x_i-\bar x)=0\). The full likelihood is thus the
normal likelihood for \(\bar x\) with variance \(\sigma^2/n\), times
\[
 c(x_1,\ldots,x_n)
 =\frac{(2\pi\sigma^2)^{-(n-1)/2}}{\sqrt n}
    \exp\!\left(-\frac{\sum_i(x_i-\bar x)^2}{2\sigma^2}\right).
\]
This positive factor is independent of \(\theta\) and common to both
models, so it multiplies both marginal likelihoods and cancels.
The sampling law of \(\bar X\) is the stated normal law by
\cref{thm:p15-normal-sums} and conditional independence.

The continuous alternative prior assigns zero mass to the singleton
\(\{\theta_0\}\). Its parameter posterior is absolutely continuous
with respect to that prior by the first proof, so the singleton remains
null. The separate point-null model instead has its own model mass.
\end{proof}
\end{annexitem}

\begin{annexitem}{annex:p15.normal-factor-sensitivity}{thm:p15-normal-factor-sensitivity}{Exact prior-scale extrema and distinct asymptotic paths}
\begin{proof}
At fixed \(z\), continuity of the formula gives \(B_{01}\to1\) as
\(r\downarrow0\). As \(r\to\infty\), the exponential factor tends
to \(e^{-z^2/2}>0\) while \(\sqrt{1+r}\to\infty\).

For the reciprocal factor, let
\[
 \ell(r)=\log B_{10}
 =-\tfrac12\log(1+r)+\frac{z^2r}{2(1+r)}.
\]
Its derivative is
\(\ell'(r)=[z^2-(1+r)]/[2(1+r)^2]\).
When \(z^2\leq1\), this is negative for \(r>0\); the supremum is
the boundary value one. When \(z^2>1\), it is positive before
\(r=z^2-1\) and negative after. At that point,
\[
 B_{10}=\frac1{|z|}\exp((z^2-1)/2).
\]
This proves the value and attainment claims.

For a proper mixing probability on positive prior variances, Tonelli
makes its marginal likelihood the corresponding mixture of the normal
marginal likelihoods. Dividing by the fixed positive null density gives
the mixture of their \(B_{10}\) values, hence no larger than their
supremum. Point mixing measures include every single normal prior, so
the supremum over mixtures is the same. For arbitrary proper priors
on \(\theta\), the likelihood at fixed \(x\) is at most
\(1/\sqrt{2\pi s^2}\), with equality at \(\theta=x\).
The ratio of this maximum to \(m_0(x)\) is \(e^{z^2/2}\), achieved by
the proper point prior at \(x\). Its dependence on the observed datum
explains why this extremization is not a prior fixed before the data.

For fixed \(\tau,\sigma>0\), write \(c=\tau^2/\sigma^2>0\).
On the first path \(r_n=cn\) and \(z_n=z\), so the fixed-\(z\)
limit already proved applies. On the second path
\(z_n^2=n\delta^2/\sigma^2\). For all sufficiently large \(n\),
\(r_n/(1+r_n)\geq1/2\), and
\[
 \log B_{01,n}
 \leq\tfrac12\log(1+cn)-\frac{n\delta^2}{4\sigma^2}\longrightarrow-\infty.
\]
For completeness, \(\log y=\int_1^y t^{-1}dt\leq2(\sqrt y-1)\)
for \(y\geq1\), so the logarithmic term divided by \(n\) tends to zero.
This proves the second limit.

At \(z=2\), direct substitution for \(r=1,3,9,99,999\) gives respectively
\[
 \sqrt2e^{-1},\quad2e^{-3/2},\quad
 \sqrt{10}e^{-9/5},\quad10e^{-99/50},\quad
 \sqrt{1000}e^{-999/500}.
\]
Dividing each by one plus itself gives the equal-prior null probabilities
in the body. These exact expressions determine the displayed roundings.
\end{proof}
\end{annexitem}

\begin{annexitem}{annex:p15.model-average-moments}{thm:p15-model-average-moments}{Mixture integration and total variance}
\begin{proof}
For every measurable \(A\), \(x\mapsto\sum_kw_k(x)R_k(x,A)\) is
measurable. For fixed \(x\), it is a finite nonnegative combination of
probability measures with weights summing to one; countable additivity
and total mass one follow. This verifies the kernel property, including
the declared exceptional-data default.

The integration identity
\(\int u\,dR=\sum_kw_k\int u\,dR_k\) holds for indicators by definition,
then for simple functions and nonnegative functions by monotone convergence.
For an integrable signed function apply it to its positive and negative
parts; on each positive-weight component these integrals are finite.
Zero-weight components are omitted.

The second-moment hypotheses imply first-moment integrability, for example
from \(|z|\leq1+z^2\). Apply the identity to \(z,z^2\). The first gives
\(a=\sum_kw_ka_k\); the second gives
\(\int z^2\,dR=\sum_kw_k(v_k+a_k^2)\).
Subtracting \(a^2\) proves the source form of the variance identity.
Expanding the squares and using \(\sum_kw_ka_k=a\), \(\sum_kw_k=1\)
gives
\(\sum_kw_k(a_k-a)^2=\sum_kw_ka_k^2-a^2\), proving the other form.
The between-model term is nonnegative and is positive if at least two
positive-weight component means differ.

The first example has component means \(-2,2\), each variance one,
and weights \(1/2\). Its mean is zero, within-model contribution one,
and between-model contribution four. The second has equal zero means
and variances 100 and zero, so its variance is 50. The normal moments
and the point-mass case are those of \cref{thm:p15-normal-normalization}.
Thus selection can either decrease or increase the reported variance
relative to averaging; the nonnegative between-model identity is the
precise universal statement.
\end{proof}
\end{annexitem}

\begin{annexitem}{annex:p15.model-label-refinement}{thm:p15-model-label-refinement}{Changing prior odds and splitting a label}
\begin{proof}
For positive compared prior masses and marginal likelihoods, apply
\cref{thm:p15-model-posterior-odds} under each prior vector.
The marginal likelihoods are unchanged, so posterior odds change by the
new prior-odds ratio divided by the old one. If a weight is zero, use the
weight formula directly instead of forming its odds.

Take compatible density versions for the duplicate marginal law, denoted
\(m_j\). If the duplicate masses are \(\alpha_{j1},\ldots,\alpha_{jr}\)
with sum \(\alpha_j\), their joint contribution to the evidence is
\(\sum_\ell\alpha_{j\ell}m_j=\alpha_jm_j\). The evidence denominator
is unchanged. Their posterior masses sum to \(w_j\), and, because their
quantity kernels are identical, their contributions to the model average
sum to \(w_jR_j\). With arbitrary density versions these equalities hold
almost everywhere by \cref{thm:p15-model-density-invariance}.

For identical predictive laws and original masses \(1/2,1/2\), weights
remain \(1/2,1/2\). Nine copies plus the second label, assigned new
uniform mass \(1/10\) each, give group mass \(9/10\), hence group
posterior \(9/10\). For a datum where \(m_A/m_B=2\), the original
posterior of \(A\) is \(2/(2+1)=2/3\); the relabeled uniform assignment
gives \(18/(18+1)=18/19\). Such a datum can occur in two proper binary
models with success probabilities .8 and .4. This example does not
assert that one normalized density is twice another at every point.
\end{proof}
\end{annexitem}

\begin{annexitem}{annex:p15.improper-model-priors}{thm:p15-improper-model-priors}{Normalization constants and prior-null restrictions}
\begin{proof}
Multiplication of the improper measure by \(c_k\) multiplies both
\(f_k(x\mid\theta)\widetilde\Pi_k(d\theta)\) and its finite positive
normalizing integral by \(c_k\). It therefore cancels within the
conditional parameter posterior. Across models the ratio instead becomes
\((c_i/c_j)Z_i/Z_j\). Tonelli applies to the sigma-finite parameter
and data measures and gives
\[
 \int Z_k(x)\nu(dx)
 =\int\left(\int f_k(x\mid\theta)\nu(dx)\right)\widetilde\Pi_k(d\theta)
 =\widetilde\Pi_k(\Theta_k)=\infty.
\]
Thus the displayed integrals cannot be normalized predictive probabilities.

If the same identified nuisance measure \(\rho\) is replaced by \(c\rho\)
in both model integrals, each finite integral is multiplied by \(c\);
the ratio cancels that scalar. The calculation says nothing about replacing
\(\rho\) by a different measure, or using unrelated parameter-specific
constants. It does not change an infinite total mass into one.

For the flat normal example, Gaussian normalization and translation give
\(\int_{\mathbb R}\phi_{\theta,1}(x)\,d\theta=1\).
The integrated value against \(c\,d\theta\) is \(c\), and the normalized
posterior density is \(\phi_{x,1}(\theta)\). For the proper uniform prior
on \([-A,A]\), the integral is positive and no larger than \(1/(2A)\);
hence it tends to zero for fixed \(x\), with no selected finite flat
constant emerging from that limit.

For the proper-prior example, set
\(R_a(\psi,\cdot)=\operatorname{Uniform}[0,1]\) for
\(\psi\ne1/2\), and \(R_a(1/2,\cdot)=\delta_a\), where \(a=0\)
or \(a=1\). Each is a measurable probability kernel because the changed
set is a Borel singleton. For every measurable \(A,B\subseteq[0,1]\),
\[
 \int_A R_a(\psi,B)\,d\psi
 =\lambda(A)\lambda(B),
\]
as the one changed \(\psi\) has Lebesgue measure zero. Both are therefore
versions of the same conditional law under the independent uniform joint
prior. If used to specify a separate null model at \(\psi=1/2\), however,
they choose \(\beta=0\) or \(\beta=1\). The binary likelihood
\(.2+.6\beta\) then gives success probability .2 or .8. This establishes
the claimed nonuniqueness without assigning a conditional probability
to a null event by division.
\end{proof}
\end{annexitem}

\begin{annexitem}{annex:p15.finite-model-evaluation}{thm:p15-finite-model-evaluation}{Finite summation, constants and exponential rescaling}
\begin{proof}
Integration against a finite discrete prior is its weighted sum. Substitute
that sum in the model-weight formula. If each model likelihood is multiplied
by the same finite positive \(c(x)\), each \(m_k\) and the total evidence
are multiplied by it. It cancels in every defined comparison. For model
factors \(c_i(x),c_j(x)\), a Bayes factor is instead multiplied by
\(c_i(x)/c_j(x)\), so parameter-independence alone is insufficient.

For the logarithmic calculation, discard zero terms, as they add nothing.
Every remaining term equals
\(e^{t_*}\exp(t_{k\ell}-t_*)\). The finite positive factor \(e^{t_*}\)
cancels between the numerator and denominator of every model weight,
proving the exact expression. No statement about machine arithmetic
is needed for this identity.

For the binomial example, full likelihood masses are
\(2(.5)(.5)=.5\) and
\(\tfrac12[2(.2)(.8)+2(.8)(.2)]=.32\).
Their ratio is \(.32/.5=.64=16/25\), and equal-prior alternative mass is
\(.32/(.32+.5)=16/41\). Dividing both likelihoods by two leaves
\(.25,.16\), which have the same ratio and weight.
At zero, the normal densities with variances one and four equal
\(1/\sqrt{2\pi}\) and \(1/(2\sqrt{2\pi})\), respectively.
Their ratio is two although their unnormalized exponential kernels
both equal one. This verifies the two cancellation examples.
\end{proof}
\end{annexitem}

\begin{annexitem}{annex:p15.normal-significance-comparison}{thm:p15-normal-significance-comparison}{The null-tail law and posterior model probabilities}
\begin{proof}
Under \(H_0\), \(Z=(X-\theta_0)/s\) is standard normal. Its density is
positive, continuous and even. Thus \(q(t)=2(1-\Phi(t))\), \(t\geq0\),
is continuous, strictly decreasing from one to zero. For each \(0<u<1\)
there is a unique \(t_u>0\) with \(q(t_u)=u\), and
\[
 \Pr_0\{q(|Z|)\leq u\}
 =\Pr_0\{|Z|\geq t_u\}=2(1-\Phi(t_u))=u.
\]
The normal law has no point masses, so endpoint conventions at \(t_u\)
do not alter this result.

At \(z=2\), the exact value is \(2(1-\Phi(2))\), approximately .0455.
It contains neither the alternative variance nor the model prior masses.
By contrast, the table's \(w_0\) values come from
\(B_{01}/(1+B_{01})\) and change with \(r\).
On the fixed-\(z\) path, \(B_{01}\to\infty\), and
\[
 w_0=\frac{\alpha_0B_{01}}{\alpha_0B_{01}+\alpha_1}\longrightarrow1
 \quad(\alpha_0,\alpha_1>0\text{ fixed}).
\]
The tail probability remains \(q(|z|)\). On the fixed nonzero-departure
path, the preceding sensitivity theorem instead gives \(B_{01}\to0\).
Both distinctions follow under their explicitly different sequences of
data and assumptions, without identifying a null tail probability with
a posterior model probability.
\end{proof}
\end{annexitem}
